% ============================================================
% SKELETON -- Part II, Section 8: Interacting Gases
% ============================================================

\section{Interacting Gases}

\paragraph{Scope of this section.}
Real particles attract at long range and repel at short range. This section
adds pair interactions to the ideal gas of Section~7, first exactly at low
density and then through two closed-form models.

\subsection{The Configurational Integral}

\paragraph{Separation of the momentum integrals.}
Pair interactions $\phi(r)$ depend on the positions only, so the momentum
integrals are those of the ideal gas. The partition function then splits
into the ideal-gas factor and a \textcolor{blue}{configurational integral}
$Q_N$.
\begin{equation}
    \boxed{Z=\frac{Q_N}{N!\,\lambda^{3N}},\qquad
    Q_N=\int d^{3N}r\;\exp\!\Big[-\beta\sum_{i<j}\phi(r_{ij})\Big].}
\end{equation}

\subsection{The Virial Expansion}

\paragraph{Equation of state at low density.}
At low density, the equation of state is a power series in $1/v$. The
\textcolor{blue}{second virial coefficient} $B_2(T)$ follows exactly from
the pair potential through the \textcolor{blue}{Mayer function}
$f(r)=e^{-\beta\phi(r)}-1$. This expansion is the rigorous link between the
interactions and the deviation from ideal behavior.
\textcolor{red}{Here, $v\equiv V/N$ denotes the volume per particle.}
% TODO: derive B_2 from Q_N to first order in the density.
\begin{equation}
    \boxed{\frac{Pv}{k_BT}=1+\frac{B_2(T)}{v}+\cdots,\qquad
    B_2(T)=-\frac{1}{2}\int d^3r\,f(r).}
\end{equation}

\paragraph{Sign of the second virial coefficient.}
The repulsive core makes $B_2$ positive, and the attraction makes $B_2$
negative. Repulsion dominates at high temperature, and attraction dominates
at low temperature.
% TODO: hard core plus square well as the worked example; Boyle temperature
%       B_2(T_B) = 0.

\subsection{The van der Waals Equation}

\paragraph{Two assumptions.}
Two assumptions turn $Q_N$ into a closed form. A hard core excludes the
volume $b$ per particle, and the attraction acts on each particle as a
uniform mean field proportional to the density, with strength $a$.
% TODO: Q_N = (V - Nb)^N exp(beta a N^2 / V); state both as assumptions.
\begin{equation}
    \boxed{P=\frac{k_BT}{v-b}-\frac{a}{v^2}.}
\end{equation}

\paragraph{Comparison with the virial expansion.}
The expansion of the van der Waals equation in $1/v$ gives a second virial
coefficient with the same two contributions as the exact $B_2$.
\begin{equation}
    \boxed{B_2^{\text{vdW}}(T)=b-\frac{a}{k_BT}.}
\end{equation}

\paragraph{Critical point and corresponding states.}
Below a critical temperature, each van der Waals isotherm has a loop, and
the middle part of the loop has $\kappa_T<0$, which violates the stability
condition of Section~6. The loop shrinks to a horizontal inflection at the
critical point, where $(\partial P/\partial v)_T=(\partial^2P/\partial
v^2)_T=0$. The reduced variables $P_r=P/P_c$, $v_r=v/v_c$ and $T_r=T/T_c$
remove the material constants.
\begin{equation}
    \boxed{v_c=3b,\quad k_BT_c=\frac{8a}{27b},\quad P_c=\frac{a}{27b^2},
    \qquad
    \left(P_r+\frac{3}{v_r^2}\right)\left(3v_r-1\right)=8T_r.}
\end{equation}
% TODO: law of corresponding states; the loop is resolved by phase
%       coexistence in Part III (Maxwell construction).

\subsection{The Dieterici Equation}

\paragraph{An alternative closed form.}
The \textcolor{blue}{Dieterici equation} replaces the mean-field attraction
by a Boltzmann factor and has the same second virial coefficient as the van
der Waals equation.
\begin{equation}
    \boxed{P=\frac{k_BT}{v-b}\exp\!\left(-\frac{a}{k_BTv}\right).}
\end{equation}

\paragraph{The critical compressibility factor as a test.}
In either model, the critical compressibility factor
\begin{equation*}
    Z_c\equiv\frac{P_cv_c}{k_BT_c}
\end{equation*}
has one value for all substances, so $Z_c$ tests both models against
experiment (Table~\ref{tab:Zc}).
% TODO: Dieterici critical point v_c = 2b, k_B T_c = a/(4b), P_c = a/(4 e^2 b^2).
\begin{table}[h]
\centering
\caption{Critical compressibility factor $Z_c$ of two model equations and
four fluids.}
\label{tab:Zc}
\begin{tabular}{lc}
    Equation or fluid & $Z_c$ \\ \hline
    van der Waals & $3/8=0.375$ \\
    Dieterici & $2/e^2\approx0.271$ \\
    Argon & 0.291 \\
    Nitrogen & 0.289 \\
    Carbon dioxide & 0.274 \\
    Water & 0.229 \\
\end{tabular}
\end{table}

\subsection{Fugacity}

\paragraph{Definition.}
The \textcolor{blue}{fugacity} $f$ keeps the logarithmic form of the
ideal-gas chemical potential of Section~7 for a real gas. The
\textcolor{blue}{fugacity coefficient} $\varphi\equiv f/P$ measures the
deviation from ideal behavior, and $\varphi\to1$ as $P\to0$.
\begin{equation}
    \boxed{\mu(T,P)=\mu^\circ(T)+k_BT\ln\frac{f}{P^\circ},\qquad
    \varphi\equiv\frac{f}{P}.}
\end{equation}

\paragraph{Fugacity from the equation of state.}
The derivative $(\partial\mu/\partial P)_T=v$ from the Gibbs--Duhem relation,
integrated from the ideal-gas limit, gives the fugacity coefficient from any
equation of state.
% FROM: Gibbs--Duhem relation (Section 5).
% TODO: derive; virial result ln(phi) = B_2 P/(k_B T) to first order;
%       van der Waals fugacity as an example.
\begin{equation}
    \boxed{\ln\varphi=\int_0^P\left(\frac{v}{k_BT}-\frac{1}{P'}\right)dP'
    \approx\frac{B_2P}{k_BT}.}
\end{equation}

\paragraph{Fugacity and phase equilibrium.}
Two phases in equilibrium have equal chemical potentials, so the two phases
have equal fugacities. Part~III uses this equality for the vapor pressure
of a real fluid.
\begin{equation}
    \boxed{f_l(T,P)=f_g(T,P)\quad\text{at coexistence.}}
\end{equation}
